\subsection{Критерий однозначной обратимости непрерывной на отрезке функции}

\textbf{Теорема.} Пусть $f(x) \in C[a, b]$.
\[ f(x) \text{ — однозначно обратима} \iff f(x) \text{ — строго монотонна} \]

\begin{tcolorbox}[title=Доказательство]
    \textbf{1. Достаточность ($\Leftarrow$)}
    Следует из \textbf{Теоремы 4} (об обратной функции):
    Если $f(x)$ непрерывна и строго монотонна, то она обратима.

    \textbf{2. Необходимость ($\Rightarrow$)}
    Дано: $f(x)$ непрерывна на $[a, b]$ и однозначно обратима (инъективна).
    Доказать: $f(x)$ строго монотонна.

    \begin{enumerate}
        \item \textbf{Проверка граничных значений.}
        $f(a) \neq f(b)$. (Если бы $f(a) = f(b)$, то значение принималось бы дважды, что противоречит обратимости).
        
        Пусть для определенности $A := f(a) < f(b) := B$.
        В этом случае докажем, что $f(x) \uparrow\uparrow$ (строго возрастает) на $[a, b]$.
        (Случай $A > B$ доказывается аналогично для убывания).

        \item \textbf{Множество значений.}
        Обозначим $Y = f([a, b])$.
        Так как функция непрерывна, $Y$ — это промежуток. Так как она обратима, каждое значение принимается один раз.
        Докажем, что $Y = [A, B]$.
        
	\begin{center}
\begin{tikzpicture}[scale=1.2]
    % Оси координат
    \draw[->] (-0.5,0) -- (6,0) node[right] {$x$};
    \draw[->] (0,-0.5) -- (0,5) node[above] {$y$};

    % Ключевые точки
    \coordinate (A_pt) at (1, 3);      % Точка (a, A)
    \coordinate (X0_pt) at (2.5, 1);   % Точка (x0, f(x0)) - "яма"
    \coordinate (C_pt) at (4, 3);      % Точка (c, A) - пересечение
    \coordinate (B_pt) at (5.5, 4.5);  % Точка (b, B)

    % График функции (кривая)
    \draw[thick, blue] (A_pt) 
        .. controls (1.2, 3) and (1.5, 1) .. (X0_pt) 
        .. controls (3.5, 1) and (3.8, 3) .. (C_pt) 
        .. controls (4.5, 3) and (5, 4.5) .. (B_pt) node[right] {$f(x)$};

    % Горизонтальный уровень A
    \draw[dashed] (0, 3) node[left] {$A$} -- (5.5, 3);
    
    % Проекции на ось X
    \draw[dashed] (A_pt) -- (1, 0) node[below] {$a$};
    \draw[dashed] (X0_pt) -- (2.5, 0) node[below] {$x_0$};
    \draw[dashed] (C_pt) -- (4, 0) node[below] {$c$};
    \draw[dashed] (B_pt) -- (5.5, 0) node[below] {$b$};

    % Выделение точек противоречия
    \fill[red] (A_pt) circle (2pt);
    \fill[red] (C_pt) circle (2pt);
    
    % Поясняющие стрелки (неравенства)
    \draw[<->] (1, -0.6) -- (2.5, -0.6) node[midway, below] {\scriptsize $x_0 \ge a$};
    \draw[<->] (2.5, -0.6) -- (4, -0.6) node[midway, below] {\scriptsize $c > x_0$};
    
    % Комментарий
    \node[red, align=left] at (2.5, 4) {Нарушение инъективности:\\ $f(a) = f(c) = A$};

\end{tikzpicture}
\end{center}

        Предположим, что $\exists y_0 = f(x_0)$ такой, что $y_0 \notin [A, B]$.
        Например, $y_0 < A$.
        Тогда имеем: $f(x_0) < f(a) < f(b)$.
        По теореме о промежуточном значении на отрезке $[x_0, b]$ функция должна принять значение $f(a)$ (так как оно лежит между $f(x_0)$ и $f(b)$).
        То есть $\exists c \in (x_0, b): f(c) = f(a)$.
        Но $c \neq a$, так как $c > x_0 \ge a$.
        Получили $f(c) = f(a)$, что противоречит однозначной обратимости.
        \textbf{Вывод:} $f([a, b]) = [A, B]$.

        \item \textbf{Доказательство монотонности (от противного).}
        Предположим, что $f(x)$ \textbf{не} является строго возрастающей.
        Это значит, что существуют точки $x_1, x_2 \in (a, b)$ такие, что:
        \[ x_1 < x_2, \quad \text{но} \quad f(x_1) \ge f(x_2) \]
        
        Так как равенство $f(x_1) = f(x_2)$ невозможно (в силу обратимости), то:
        \[ f(x_1) > f(x_2) \]
        
        Рассмотрим значения:
        \begin{itemize}
            \item $f(a) = A$
            \item $f(b) = B$
            \item $f(x_1) > f(x_2)$
        \end{itemize}
        
        Возможны ситуации (см. рисунок "Зигзаг"):
        Поскольку $f(x_2) < f(x_1)$ и $f(x_2) \ge A$, а $f(b) = B \ge f(x_1)$ (так как мы доказали в п.2, что значения не выходят за $[A, B]$? \textit{Или даже если не доказали}).
        
        Рассмотрим отрезок $[x_2, b]$.
        На левом конце значение $f(x_2)$. На правом $f(b) = B$.
        Заметим, что $f(x_2) < f(x_1)$.
        Также, так как $f(x_1)$ — значение функции, то $f(x_1) \in [A, B]$, значит $f(x_1) \le f(b)$.
        Если $f(x_1) < f(b)$, то число $f(x_1)$ лежит между $f(x_2)$ и $f(b)$.
        
        По теореме о промежуточном значении:
        \[ \exists x' \in [x_2, b]: \;\; f(x') = f(x_1) \]
        
        Но $x' \ge x_2 > x_1$, то есть $x' \neq x_1$.
        Мы нашли две разные точки $x_1$ и $x'$, в которых функция принимает одинаковое значение.
        \[ f(x_1) = f(x') \implies \text{Противоречие с обратимостью!} \]
        
        (Если бы оказалось, что $f(x_1) > f(b)$, мы бы искали противоречие на отрезке $[a, x_1]$, найдя там точку, равную $f(b)$).
        
        \begin{center}
        \begin{tikzpicture}[scale=0.8]
            \draw[->] (-0.5,0) -- (5,0) node[right] {$x$};
            \draw[->] (0,-0.5) -- (0,4) node[above] {$y$};
            
            % Точки
            \coordinate (A) at (0.5, 1);
            \coordinate (X1) at (2, 3);
            \coordinate (X2) at (3, 1.5);
            \coordinate (B) at (4.5, 3.5);
            
            % График (зигзаг)
            \draw[thick] (A) -- (X1) -- (X2) -- (B);
            
            % Подписи осей
            \draw[dashed] (A) -- (0.5, 0) node[below] {$a$};
            \draw[dashed] (B) -- (4.5, 0) node[below] {$b$};
            \draw[dashed] (X1) -- (2, 0) node[below] {$x_1$};
            \draw[dashed] (X2) -- (3, 0) node[below] {$x_2$};
            
            \draw[dashed] (A) -- (0, 1) node[left] {$A$};
            \draw[dashed] (B) -- (0, 3.5) node[left] {$B$};
            
            % Иллюстрация противоречия
            \draw[dashed, red] (0, 3) node[left] {$f(x_1)$} -- (4, 3);
            \fill[red] (X1) circle (2pt);
            \fill[red] (3.8, 3) circle (2pt) node[above left] {$x'$};
            
            \node at (2.5, -1) {Нарушение монотонности $\to$ нарушение инъективности};
        \end{tikzpicture}
        \end{center}
    \end{enumerate}
    \textbf{Итог:} Предположение неверно, функция обязана быть строго монотонной.
\end{tcolorbox}
